problem:  
Let \((M^6, J, g)\) be a compact, simply connected, **strict nearly Kähler** manifold with SU(3)-structure forms \((\omega, \psi_+, \psi_-)\).  
The canonical Hermitian connection \(\nabla^c\) on \(M\) satisfies \(\nabla^c T = 0\), where \(T = \frac{1}{2} d\omega\) is its totally skew-symmetric torsion.  
It follows that \({\rm Hol}(\nabla^c) \subseteq SU(3)\), and for all **known** compact examples (both homogeneous and inhomogeneous, e.g. those found by Foscolo–Haskins), one has full holonomy \({\rm Hol}(\nabla^c) = SU(3)\).  

Now, consider a **smooth one-parameter family of strict nearly Kähler SU(3)-structures** \((\omega_t, \psi_{+,t}, \psi_{-,t})\) on a fixed compact \(6\)-manifold \(M\), deforming within the strict nearly Kähler condition:  
\[
d\omega_t = 3\,\psi_{+,t}, \quad d\psi_{-,t} = -2\,\omega_t^2.
\]  
Denote by \(\nabla^c_t\) the corresponding family of canonical Hermitian connections, with torsions \(T_t = \frac{1}{2}d\omega_t\).  

**Problem:**  
Is the holonomy type of the canonical connection \({\rm Hol}(\nabla^c_t)\) **analytically stable** under nearly Kähler deformations? That is:  

1. If \({\rm Hol}(\nabla^c_0) = SU(3)\), must \({\rm Hol}(\nabla^c_t)\) remain \(SU(3)\) for small \(t\neq 0\)?  
2. Could a nearly Kähler deformation trigger a *reduction* of the canonical connection’s holonomy (e.g. to a proper subgroup of \(SU(3)\)) through analytic degeneration of the torsion tensor, even while preserving strictness?  
3. If such holonomy jump phenomena are impossible, can one express a formal *holonomy rigidity condition* (in terms of higher-order Laplace-type operators on intrinsic torsion forms) ensuring that \(\nabla^c\)-parallel torsion forces the holonomy to remain maximal under smooth nearly Kähler deformations?

Why is it a "good" problem:  
This problem refines the holonomy question to an *analytic stability* problem under deformations, thus moving beyond existing structure theorems. It does not ask whether smaller holonomy exists (which is known to be false in the compact strict case), but whether *holonomy reduction could occur through continuous deformation*. Such a deformation‑stability problem is currently unaddressed in the literature and would enrich our understanding of how the SU(3) structure interacts with its canonical connection when varying within the nearly Kähler class.  

It connects three deep mathematical themes:  
- **Holonomy theory with torsion** (exploring how \({\rm Hol}(\nabla^c)\) behaves under smooth deformations).  
- **Analytic deformation theory of special geometries**, using elliptic operators governing infinitesimal nearly Kähler deformations.  
- **Geometric analysis of G₂ cones**, because deformations of the NK structure correspond to deformations of the associated G₂ cone metric (and holonomy of the cone controls the SU(3)‑holonomy base).

The question is short, precise, and meaningful but leads directly into unknown analytic territory, potentially revealing a rigidity or bifurcation phenomenon linking SU(3)‑geometry, G₂ holonomy, and harmonic analysis on manifolds with torsion—hence a “good problem” in the deepest sense.